%% Supplementary material.  Compiled on its own (latexmk -pdf supplementary.tex);
%% the figure sheets and table bodies in supplementary/ are made by
%% scripts/fig_crack3d_supp.py, fig_crackdir_supp.py, fig_tau_geodesic.py and
%% supp_tables.py.
\documentclass[preprint,12pt]{elsarticle}
\usepackage{amsmath}
\usepackage{amssymb}
\usepackage{siunitx}
\usepackage{graphicx}
\usepackage{booktabs}
\usepackage{longtable}
\usepackage{array}
\usepackage{url}
\renewcommand{\thetable}{S\arabic{table}}
\renewcommand{\thefigure}{S\arabic{figure}}
\renewcommand{\thesection}{S\arabic{section}}

\begin{document}
\begin{frontmatter}
  \title{Supplementary material to\\Grain-scale kinematics, ice-bond failure
    and damage evolution in ice-bonded granular media under in-situ uniaxial
    compression}
  \author{Amirali Khosrozadeh, Yannik Sinnwell, Swantje Pietsch-Braune, Kai
    Nikolaus, Sergiy Antonyuk, Stefan Heinrich}
\end{frontmatter}

The analysis code and the tables reproduced here are at
\url{https://github.com/AmiraliKhosrozadeh/ice-bonded-granular-microCT}.
Specimen labels, load-step numbering and the axial convention (punch at the
top) are those of the paper.

\section{Method settings and provenance}
Table~\ref{tab:settings} lists, for each analysis of the paper, what it
takes in, what it returns, the software that performs it and the settings
that fix the result.  Every entry not attributed to a package is the
authors' own code.

  {\footnotesize
   \setlength{\tabcolsep}{4pt}
   \renewcommand{\arraystretch}{1.15}
   \begin{longtable}{p{2.4cm}p{2.6cm}p{2.4cm}>{\raggedright\arraybackslash}p{5.6cm}}
      \caption{Inputs, outputs, software and settings of each analysis.}
      \label{tab:settings} \\
      \toprule
      Analysis & Input & Output & Software and settings \\
      \midrule
      \endfirsthead
      \toprule
      Analysis & Input & Output & Software and settings \\
      \midrule
      \endhead
      Phase classification & filtered greyscale volume & phase map (air, ice, grain) & own code; $3^3$ median; valley thresholds per scan for beads, three-class Otsu for sand \\
      Bead labels & phase map & grain labels & own code; marker-controlled watershed on the distance transform; shape gate on fill, radius and roughness \\
      Sand labels & greyscale volume & grain labels & own 3D U-Net, border--core, trained on synthetic packings; selected against watershed, a uniform-intensity U-Net, ParticleSeg3D and Cellpose-SAM (Table~2 of the paper) \\
      Bead DDIC & greyscale scans $N$, $N{+}1$; labels of $N$ & per-grain $\Phi$; translation used & SPAM~0.9 \texttt{spam-ddic}; full affine, label dilation 2, margin 15 voxels, 300 iterations, two-level multiscale, updated gradient, per-grain seed; converged at return status~2, usable at~1; grey scales mapped on a copy where they differed \\
      Sand tracking & labels of both scans; column lengths & matched grain pairs & own code; affine predictor from column length, refinement over the 12 nearest accepted matches, half-spacing tolerance, volume ratio 0.6--1.6, reciprocity required \\
      Strain and rotation & accepted translations & $\mathbf{G}$, $\mathbf{E}$, $E_\mathrm{eq}$, rotation & own code; least squares with intercept over 12 (beads) or 16 (sand) nearest neighbours; Green--Lagrange; Eq.~(7) of the paper; composed $1\!\to\!N$ fields re-fitted at scan-1 positions; rigid column motion removed \\
      Bond census & phase map, labels, matched pairs & throats, $\bar z$, neck radius, survival, coverage pair, class & own code; lens footprint $0.40R$; neck on the plane through the narrowest free point; survival classes at 12 voxels; coverage in a 1--3 voxel shell on a $3\times$ upsampled grid; $C_\mathrm{HI} \in \{0.5,0.7,0.8,0.9\}$, reference 0.8; answerable at gap $\geq 7$ voxels \\
      Breakage clustering & broken-bond midpoints & cluster size, RMS, $p$ & scikit-learn DBSCAN, min.\ samples 5, clusters $\geq 6$; neighbourhood one grain diameter, sweep 0.6--2.0; 400 resamplings of the bonds present; Benjamini--Hochberg across specimens \\
      Core structure & phase map & $\varepsilon_\mathrm{c}$, $\phi_\mathrm{ice}$, $\rho$, $S_V$, spanning air, $k/k_0$ & own code; core radius $0.85$ of the 98th-percentile grain-cloud radius, height band 0.20--0.80 of the column, re-centred per scan; $S_V$ by voxel-face count; 26-connectivity on the $2\times$-binned map, spanning between the terminal 5\,\% bands; Carman form with $k_K = 5$ \\
      Crack classes & phase map, labels & void, surface and body crack, separation & own code; bead envelope (closing $1.25R$, holes filled, 4-voxel skin, one bead diameter clear of each platen); void ceiling \SI{2}{\milli\meter\cubed}; crack width $\leq$ \SI{0.25}{\milli\meter}; skin 20 voxels; bodies $< 100$ voxels dropped; change from the unloaded scan reported \\
      Ice tortuosity & ice phase of interior cubes & $D_{\mathrm{eff},z}$, $\tau_\mathrm{ice}$ & pumapy~3.2.2 \texttt{compute\_continuum\_tortuosity}; $2\times$-binned map; $128^3$ voxels (\SI{6.3}{\milli\meter}); three cubes at 0.28--0.72 of the column on the axis, $\geq 98\%$ inside the specimen; $D_0 = 1$ in ice, 0 elsewhere; $c = 0/1$ on the axial faces, periodic sides; conjugate gradient, tolerance $10^{-6}$, $2\times10^5$ iterations; accepted if $D_{\mathrm{eff},z} > 0$ and $\tau \geq 1$; mean over accepted cubes \\
      Ice-path map & phase map & detour of the shortest ice path from the support forced by the crack air & own code, scikit-image \texttt{MCP\_Geometric}, unit cost; grain envelope (per-slice closing $1.25R$, holes filled, 2-voxel skin; sand pulled in \SI{1}{\milli\meter}); crack air = air inside the envelope removed by a \SI{0.25}{\milli\meter} opening; largest 26-connected body of ice plus crack air, ice taken inside it; seeds on the first ice voxel of every axial line within \SI{3}{\milli\meter} of the support, the same for both paths; drawn where the path through the ice exceeds the path through the healed ice by more than \SI{0.3}{\milli\meter}, opened once, bodies $\geq$ \SI{0.5}{\milli\meter\cubed}; $2\times$-binned maps for beads, full resolution for sand \\
      \bottomrule
   \end{longtable}}

\section{Void and crack classes in every bead specimen}
Figures~\ref{fig:crack_glass} and \ref{fig:crack_alumina} give the classes
of \S2.5.2 of the paper (void, surface crack, body crack inside the bead
envelope) at every scan of every bead specimen, in the rendering of Fig.~11
of the paper.  Table~\ref{tab:crack} gives the volumes.

\begin{figure}[p]
  \centering
  \includegraphics[width=\linewidth,height=0.9\textheight,keepaspectratio]{supplementary/crack3d_glass.png}
  \caption{Void and crack classes at every scan of the glass specimens.}
  \label{fig:crack_glass}
\end{figure}
\begin{figure}[p]
  \centering
  \includegraphics[width=\linewidth,height=0.9\textheight,keepaspectratio]{supplementary/crack3d_alumina.png}
  \caption{Void and crack classes at every scan of the $\gamma$-alumina specimens.}
  \label{fig:crack_alumina}
\end{figure}

\begin{table}[p]
  \centering
  \caption{Volumes of the classes inside the bead envelope at every scan.
    The unloaded scan is the baseline of each specimen; the floor of a
    change is set by A4's first increment (see the paper).}
  \label{tab:crack}
  \footnotesize
  \begin{tabular}{lcccccc}
\toprule
Specimen & scan & envelope (mm$^3$) & void & surface crack & body crack & separation (mm$^3$) \\
\midrule
G1 & 1 & 1356 & 6.8 & 3.0 & 0.1 & 2.6 \\
G1 & 2 & 1358 & 7.8 & 61.1 & 44.9 & 52.8 \\
G1 & 3 & 1383 & 13.7 & 49.0 & 55.9 & 18.9 \\
G2 & 1 & 1148 & 9.6 & 0.0 & 0.0 & 0.0 \\
G2 & 2 & 1195 & 9.8 & 52.7 & 30.4 & 34.2 \\
G3 & 1 & 1514 & 34.2 & 0.0 & 0.0 & 0.0 \\
G3 & 2 & 1486 & 12.8 & 6.0 & 1.4 & 3.1 \\
G3 & 3 & 1614 & 3.8 & 49.3 & 22.9 & 32.0 \\
G3 & 4 & 1344 & 11.7 & 92.4 & 49.3 & 63.6 \\
G4 & 1 & 1242 & 6.2 & 0.0 & 0.0 & 0.0 \\
G4 & 2 & 1220 & 16.0 & 30.5 & 20.2 & 14.6 \\
G5 & 1 & 1146 & 10.5 & 7.6 & 1.0 & 4.0 \\
G5 & 2 & 1180 & 8.3 & 80.7 & 26.0 & 43.7 \\
A1 & 1 & 1358 & 55.7 & 12.2 & 26.8 & 25.3 \\
A1 & 2 & 1309 & 61.7 & 9.7 & 20.3 & 17.5 \\
A1 & 3 & 1314 & 45.7 & 33.2 & 36.8 & 33.6 \\
A2 & 1 & 1388 & 46.5 & 21.7 & 39.8 & 37.2 \\
A2 & 2 & 1477 & 22.6 & 103.0 & 75.5 & 83.8 \\
A3 & 1 & 1834 & 75.9 & 18.6 & 25.2 & 16.3 \\
A3 & 2 & 1279 & 37.6 & 37.7 & 41.2 & 23.7 \\
A4 & 1 & 1447 & 51.2 & 28.5 & 38.4 & 45.4 \\
A4 & 2 & 1500 & 27.7 & 72.0 & 57.9 & 79.6 \\
A4 & 3 & 1528 & 27.3 & 102.2 & 65.0 & 86.1 \\
\bottomrule
\end{tabular}

\end{table}

\section{Crack orientation in every bead specimen}
The pipeline's body-crack components with their fitted planes, in the
rendering of Fig.~12 of the paper, for every loaded scan
(Figs.~\ref{fig:dir_glass} and \ref{fig:dir_alumina}).  These renders use
the pipeline's per-scan specimen mask, whose boundary artefact the paper
describes in \S2.5.2; the fitted orientations are read as the principal
direction of a damage zone, not as fracture planes.

\begin{figure}[p]
  \centering
  \includegraphics[width=\linewidth,height=0.9\textheight,keepaspectratio]{supplementary/crackdir_glass.png}
  \caption{Crack orientation classes of the glass specimens.}
  \label{fig:dir_glass}
\end{figure}
\begin{figure}[p]
  \centering
  \includegraphics[width=\linewidth,height=0.9\textheight,keepaspectratio]{supplementary/crackdir_alumina.png}
  \caption{Crack orientation classes of the $\gamma$-alumina specimens.}
  \label{fig:dir_alumina}
\end{figure}

\section{Ice-path maps of every specimen}
Figures~\ref{fig:tau_glass}--\ref{fig:tau_sand} give the map of Fig.~18 of
the paper for every specimen.  For every ice voxel the shortest path
through the ice from the support face is computed with a cost inversely
proportional to the local ice fraction in a \SI{0.5}{\milli\meter} window
(normalised by the part of the window inside the grain envelope, so the
lateral surface does not read as porous).  The support face is treated as
an equal-potential boundary, as in the diffusion solve: every ice voxel
that is the first of its axial line from the support, within
\SI{3}{\milli\meter} of it, seeds the path, and the axial distance is
measured from that seed surface.  The ratio of path to axial distance is
divided by its median at the same height in the unloaded column of the same
specimen, and ice above the 99.5th percentile of the unloaded column at the
same height (bodies of at least \SI{0.5}{\milli\meter\cubed}) is drawn
solid and coloured, the remaining ice as a faint shell, so an unloaded
column is blank by construction and a loaded one shows only what exceeds
the unloaded range.  The map is computed inside the grain envelope on the
largest 26-connected ice body, on the $2\times$-binned phase maps for the
beads and the full-resolution phase maps for the sand, whose envelope is
pulled in by \SI{1}{\milli\meter} to leave out the ice-rich skin the sand
packs against the tube wall; the \SI{6}{\milli\meter} next to the support
face are cut, and a line separates the unloaded scan from the loaded ones.
The paper shows G1, G3, A4 and S3; the share of the ice above the unloaded
range is given in Table~\ref{tab:taumap}.  In the glass it follows the crack
volume of Table~\ref{tab:crack}, G1 and G3 at their later steps far above
the rest, G2 and G5 a small share each and G4 nothing.  In the
$\gamma$-alumina A1 and A3 show nothing, and this is not a matter of the
threshold, since the distribution of their loaded fields coincides with
the unloaded one at every percentile (95th percentile $1.19$ against
$1.19$ in A1, $1.15$ against $1.15$ in A3), so any threshold that coloured
them would colour their unloaded columns as much; the PuMA tortuosity of
Table~
ef{tab:tau} and the crack classes of Table~
ef{tab:crack} are
unchanged in the same two specimens.  A2 shows a single body near the support
despite the largest crack-volume gain of the four, and A4 the two bead
layers under the punch at both steps, present already at its first
increment (shortening $0.01$), which is the punch bedding into the top
layer rather than an interior crack; two thirds or more of the cracked
throats of A2 and A4 part at the grain face, so the ice network stays
whole.  In the sand S1 shows nothing (its second scan is at zero
shortening), S2 the outward flow of the head at its third step, and S3 the
loaded third of the column with a tongue down the axis.

\begin{table}[htbp]
  \centering
  \caption{Share of the ice (inside the envelope, above the
    \SI{6}{\milli\meter} cut) whose relative path exceeds the unloaded
    column's range at the same height, per scan, in per cent.}
  \label{tab:taumap}
  \footnotesize
  \begin{tabular}{lrrrr}
\toprule
Specimen & unloaded & load step 2 & load step 3 & load step 4 \\
\midrule
G1 & 0.00 & 20.48 & 7.29 & -- \\
G2 & 0.00 & 2.44 & -- & -- \\
G3 & 0.00 & 0.00 & 5.76 & 5.00 \\
G4 & 0.00 & 0.00 & -- & -- \\
G5 & 0.00 & 1.41 & -- & -- \\
A1 & 0.00 & 0.00 & 0.00 & -- \\
A2 & 0.00 & 0.26 & -- & -- \\
A3 & 0.00 & 0.00 & -- & -- \\
A4 & 0.00 & 1.58 & 2.39 & -- \\
S1 & 0.00 & 0.00 & -- & -- \\
S2 & 0.00 & 0.00 & 0.35 & -- \\
S3 & 0.00 & 4.48 & -- & -- \\
\bottomrule
\end{tabular}

\end{table}

\begin{figure}[p]
  \centering
  \includegraphics[width=\linewidth,height=0.9\textheight,keepaspectratio]{supplementary/tau_geodesic_glass.png}
  \caption{Ice-path maps of the glass specimens.}
  \label{fig:tau_glass}
\end{figure}
\begin{figure}[p]
  \centering
  \includegraphics[width=\linewidth,height=0.9\textheight,keepaspectratio]{supplementary/tau_geodesic_alumina.png}
  \caption{Ice-path maps of the $\gamma$-alumina specimens.}
  \label{fig:tau_alumina}
\end{figure}
\begin{figure}[p]
  \centering
  \includegraphics[width=\linewidth,height=0.9\textheight,keepaspectratio]{supplementary/tau_geodesic_sand.png}
  \caption{Ice-path maps of the sand specimens.}
  \label{fig:tau_sand}
\end{figure}

\section{Per-specimen tables}

\begin{table}[p]
  \centering
  \caption{Ice tortuosity of every scan from the PuMA continuum diffusion
    solve, mean and standard deviation over the accepted $128^3$-voxel
    cubes, with the ice fraction of the cubes and the number accepted of
    those attempted.}
  \label{tab:tau}
  \footnotesize
  \begin{tabular}{lcccccc}
\toprule
Specimen & scan & shortening & $\phi_\mathrm{ice}$ (cubes) & $\tau_\mathrm{ice}$ & sd & cubes \\
\midrule
G1 & 1 & 0.000 & 0.399 & 1.609 & 0.031 & 3/3 \\
G1 & 2 & 0.161 & 0.358 & 1.888 & 0.000 & 1/1 \\
G1 & 3 & 0.224 & 0.318 & 3.737 & 2.001 & 2/2 \\
G2 & 1 & 0.000 & 0.412 & 1.622 & 0.006 & 2/2 \\
G2 & 2 & 0.206 & 0.406 & 1.817 & 0.050 & 2/2 \\
G3 & 1 & 0.000 & 0.411 & 1.610 & 0.016 & 3/3 \\
G3 & 2 & 0.045 & 0.416 & 1.605 & 0.003 & 3/3 \\
G3 & 3 & 0.102 & 0.421 & 1.620 & 0.027 & 2/2 \\
G3 & 4 & 0.302 & 0.329 & 1.840 & 0.047 & 3/3 \\
G4 & 1 & 0.000 & 0.469 & 1.520 & 0.182 & 3/3 \\
G4 & 2 & 0.219 & 0.377 & 1.889 & 0.198 & 3/3 \\
G5 & 1 & 0.000 & 0.409 & 1.608 & 0.017 & 2/2 \\
G5 & 2 & 0.139 & 0.359 & 1.894 & 0.141 & 2/2 \\
A1 & 1 & 0.000 & 0.367 & 1.762 & 0.032 & 3/3 \\
A1 & 2 & 0.061 & 0.360 & 1.818 & 0.002 & 2/2 \\
A1 & 3 & 0.119 & 0.368 & 1.765 & 0.022 & 2/2 \\
A2 & 1 & 0.000 & 0.323 & 1.950 & 0.088 & 3/3 \\
A2 & 2 & 0.100 & 0.368 & 1.949 & 0.000 & 1/1 \\
A3 & 1 & 0.000 & 0.368 & 1.895 & 0.110 & 3/3 \\
A3 & 2 & 0.367 & 0.373 & 2.011 & 0.039 & 2/2 \\
A4 & 1 & 0.000 & 0.318 & 1.920 & 0.051 & 3/3 \\
A4 & 2 & 0.005 & 0.334 & 1.864 & 0.000 & 1/1 \\
A4 & 3 & 0.077 & 0.332 & 2.105 & 0.151 & 2/2 \\
S1 & 1 & 0.000 & 0.520 & 2.330 & 0.003 & 2/2 \\
S1 & 2 & 0.000 & 0.528 & 2.273 & 0.013 & 2/2 \\
S2 & 1 & 0.000 & 0.537 & 2.189 & 0.108 & 3/3 \\
S2 & 2 & 0.048 & 0.549 & 2.168 & 0.155 & 3/3 \\
S2 & 3 & 0.248 & 0.536 & 2.353 & 0.126 & 2/2 \\
S3 & 1 & 0.000 & 0.480 & 2.672 & 0.114 & 3/3 \\
S3 & 2 & 0.086 & 0.434 & 3.235 & 0.900 & 3/3 \\
\bottomrule
\end{tabular}

\end{table}

\begin{table}[p]
  \centering
  \caption{Kinematics of every composed field, on the grains carried through
    the whole chain: median relative displacement with the rigid motion of
    the column removed, the rigid rotation of the column, the median and 95th
    percentile of the local rotation, and of the equivalent strain.}
  \label{tab:kin}
  \footnotesize
  \setlength{\tabcolsep}{3pt}
  \begin{tabular}{lcccccccc}
\toprule
Specimen & scans & grains & $|\mathbf{u}-\mathbf{u}_\mathrm{rigid}|$ (mm) & column rotation ($^\circ$) & $\theta$ median & $\theta$ p95 & $E_\mathrm{eq}$ median & p95 \\
\midrule
G1 & 1$\to$2 & 264/364 & 1.18 & 0.8 & 8.8 & 15.3 & 0.315 & 0.544 \\
G1 & 1$\to$3 & 264/364 & 1.44 & 1.4 & 10.0 & 17.2 & 0.375 & 0.695 \\
G2 & 1$\to$2 & 302/302 & 1.65 & 1.7 & 9.0 & 16.6 & 0.403 & 0.630 \\
G3 & 1$\to$2 & 254/375 & 0.34 & 6.3 & 5.4 & 17.0 & 0.006 & 0.304 \\
G3 & 1$\to$3 & 254/375 & 1.01 & 8.5 & 9.1 & 19.6 & 0.230 & 0.555 \\
G3 & 1$\to$4 & 254/375 & 1.35 & 11.7 & 10.9 & 25.5 & 0.312 & 0.762 \\
G4 & 1$\to$2 & 354/354 & 1.54 & 0.9 & 9.3 & 19.1 & 0.323 & 0.546 \\
G5 & 1$\to$2 & 278/278 & 1.30 & 1.6 & 9.1 & 18.1 & 0.379 & 0.573 \\
A1 & 1$\to$2 & 223/290 & 0.70 & 1.2 & 10.3 & 21.5 & 0.235 & 0.426 \\
A1 & 1$\to$3 & 223/290 & 1.20 & 2.0 & 14.1 & 26.5 & 0.370 & 0.866 \\
A2 & 1$\to$2 & 298/298 & 0.99 & 1.0 & 10.1 & 22.6 & 0.320 & 0.700 \\
A3 & 1$\to$2 & 2390/2390 & 0.57 & 0.2 & 9.7 & 36.5 & 0.271 & 0.846 \\
A4 & 1$\to$2 & 291/330 & 0.04 & 2.2 & 0.4 & 6.6 & 0.002 & 0.190 \\
A4 & 1$\to$3 & 291/330 & 0.97 & 2.3 & 10.3 & 22.6 & 0.277 & 0.533 \\
\bottomrule
\end{tabular}

\end{table}

\begin{table}[p]
  \centering
  \caption{Failure-mode census of every bead specimen at each coverage
    threshold, cohesive/adhesive, with the number of answerable throats, the
    median gap of those throats and the median share of the throat crack
    lying against a grain face.}
  \label{tab:fm}
  \footnotesize
  \begin{tabular}{lccccccc}
\toprule
Specimen & answerable & gap median (\si{\micro\meter}) & face share median & $C_\mathrm{HI}=0.5$ & 0.7 & 0.8 & 0.9 \\
\midrule
G1 & 60 & 208 & 0.19 & 53/7 & 30/30 & 16/44 & 8/52 \\
G2 & 14 & 180 & 0.14 & 7/7 & 4/10 & 2/12 & 0/14 \\
G3 & 30 & 215 & 0.17 & 2/28 & 0/30 & 0/30 & 0/30 \\
G4 & 12 & 246 & 0.08 & 6/6 & 4/8 & 2/10 & 0/12 \\
G5 & 28 & 220 & 0.37 & 8/20 & 2/26 & 1/27 & 1/27 \\
A1 & 27 & 202 & 0.19 & 20/7 & 15/12 & 12/15 & 9/18 \\
A2 & 50 & 213 & 0.30 & 27/23 & 19/31 & 17/33 & 10/40 \\
A3 & 178 & 215 & 0.42 & 80/98 & 20/158 & 5/173 & 0/178 \\
A4 & 32 & 210 & 0.34 & 19/13 & 12/20 & 8/24 & 3/29 \\
all & 431 & 212 & 0.33 & 222/209 & 106/325 & 63/368 & 31/400 \\
\bottomrule
\end{tabular}

\end{table}

\begin{table}[p]
  \centering
  \caption{Breakage-event clustering sweep of every bead specimen: at each
    neighbourhood radius (in grain diameters) the number of DBSCAN clusters,
    the largest cluster and the smallest RMS out-of-plane distance of a
    cluster, each against the median of 400 resamplings of the bonds present
    and the resampling $p$.  The pre-specified radius is $1.0$.}
  \label{tab:cluster}
  \scriptsize
  \renewcommand{\arraystretch}{0.82}
  \begin{tabular}{lcccc}
\toprule
Specimen & $R_\mathrm{nb}/d$ & clusters / null / $p$ & largest / null / $p$ & RMS/$R$ / null / $p$ \\
\midrule
G1 & 2.00 & 1 / 1 / 1.000 & 115 / 115 / 0.725 & 2.01 / 1.88 / 0.930 \\
G1 & 1.25 & 3 / 2 / 0.443 & 51 / 64 / 0.850 & 0.35 / 0.96 / 0.220 \\
G1 & 1.00 & 4 / 3 / 0.445 & 40 / 44 / 0.688 & 0.35 / 0.50 / 0.260 \\
G1 & 0.80 & 3 / 4 / 0.853 & 11 / 16 / 0.890 & 0.40 / 0.30 / 0.775 \\
G1 & 0.60 & 0 / 1 / 1.000 & 0 / 7 / 1.000 & -- \\
G1 & 1.50 & 2 / 1 / 0.350 & 63 / 107 / 0.955 & 1.58 / 1.74 / 0.347 \\
G2 & 2.00 & 1 / 1 / 1.000 & 115 / 115 / 0.963 & 2.19 / 2.07 / 0.963 \\
G2 & 1.50 & 1 / 1 / 1.000 & 112 / 109 / 0.115 & 2.14 / 2.00 / 0.960 \\
G2 & 1.25 & 2 / 3 / 0.920 & 81 / 71 / 0.225 & 0.79 / 0.56 / 0.762 \\
G2 & 1.00 & 7 / 4 / 0.050 & 20 / 29 / 0.892 & 0.21 / 0.34 / 0.105 \\
G2 & 0.80 & 3 / 3 / 0.507 & 10 / 10 / 0.573 & 0.11 / 0.28 / 0.008 \\
G2 & 0.60 & 0 / 0 / 1.000 & 0 / 0 / 1.000 & -- \\
G3 & 0.60 & 0 / 0 / 1.000 & 0 / 0 / 1.000 & -- \\
G3 & 0.80 & 1 / 0 / 0.003 & 12 / 0 / 0.000 & 0.42 / 0.27 / 1.000 \\
G3 & 2.00 & 1 / 1 / 1.000 & 33 / 40 / 0.938 & 1.21 / 2.05 / 0.058 \\
G3 & 1.25 & 1 / 1 / 0.675 & 27 / 7 / 0.000 & 1.17 / 0.53 / 0.996 \\
G3 & 1.50 & 1 / 2 / 0.995 & 28 / 16 / 0.080 & 1.23 / 0.77 / 0.809 \\
G3 & 1.00 & 2 / 0 / 0.003 & 16 / 0 / 0.000 & 0.62 / 0.41 / 0.864 \\
G4 & 2.00 & 1 / 1 / 1.000 & 122 / 122 / 0.568 & 2.43 / 2.41 / 0.562 \\
G4 & 1.50 & 1 / 3 / 1.000 & 114 / 58 / 0.000 & 2.43 / 0.58 / 1.000 \\
G4 & 1.25 & 4 / 4 / 0.700 & 63 / 52 / 0.037 & 0.34 / 0.34 / 0.500 \\
G4 & 0.80 & 3 / 2 / 0.175 & 13 / 9 / 0.220 & 0.29 / 0.35 / 0.277 \\
G4 & 0.60 & 1 / 0 / 0.225 & 7 / 0 / 0.120 & 0.25 / 0.32 / 0.144 \\
G4 & 1.00 & 5 / 3 / 0.205 & 20 / 38 / 0.922 & 0.28 / 0.32 / 0.360 \\
G5 & 1.00 & 4 / 3 / 0.432 & 58 / 54 / 0.403 & 0.29 / 0.36 / 0.310 \\
G5 & 2.00 & 1 / 1 / 1.000 & 118 / 118 / 0.843 & 1.86 / 2.11 / 0.007 \\
G5 & 1.50 & 1 / 2 / 1.000 & 115 / 105 / 0.113 & 1.77 / 1.09 / 0.540 \\
G5 & 1.25 & 4 / 3 / 0.168 & 70 / 78 / 0.855 & 0.29 / 0.49 / 0.135 \\
G5 & 0.80 & 5 / 2 / 0.040 & 15 / 14 / 0.472 & 0.22 / 0.38 / 0.093 \\
G5 & 0.60 & 0 / 0 / 1.000 & 0 / 0 / 1.000 & -- \\
A1 & 0.80 & 0 / 1 / 1.000 & 0 / 6 / 1.000 & -- \\
A1 & 1.00 & 5 / 4 / 0.263 & 14 / 22 / 0.858 & 0.30 / 0.36 / 0.307 \\
A1 & 1.25 & 6 / 2 / 0.005 & 22 / 66 / 1.000 & 0.38 / 0.99 / 0.193 \\
A1 & 1.50 & 1 / 1 / 1.000 & 109 / 105 / 0.065 & 2.48 / 2.09 / 1.000 \\
A1 & 2.00 & 1 / 1 / 1.000 & 111 / 111 / 0.802 & 2.46 / 2.21 / 0.998 \\
A1 & 0.60 & 0 / 0 / 1.000 & 0 / 0 / 1.000 & -- \\
A2 & 1.00 & 9 / 3 / 0.000 & 20 / 45 / 0.995 & 0.21 / 0.39 / 0.155 \\
A2 & 0.60 & 0 / 1 / 1.000 & 0 / 6 / 1.000 & -- \\
A2 & 0.80 & 2 / 2 / 0.750 & 15 / 11 / 0.242 & 0.37 / 0.33 / 0.606 \\
A2 & 2.00 & 1 / 1 / 1.000 & 121 / 121 / 0.532 & 2.67 / 2.29 / 1.000 \\
A2 & 1.25 & 5 / 1 / 0.007 & 67 / 96 / 0.897 & 0.21 / 1.67 / 0.065 \\
A2 & 1.50 & 2 / 1 / 0.058 & 88 / 114 / 0.995 & 0.93 / 2.13 / 0.052 \\
A3 & 2.00 & 1 / 1 / 1.000 & 2052 / 2052 / 0.517 & 4.91 / 4.91 / 0.542 \\
A3 & 0.60 & 51 / 33 / 0.000 & 22 / 19 / 0.320 & 0.25 / 0.22 / 0.830 \\
A3 & 1.25 & 1 / 1 / 1.000 & 2032 / 2033 / 0.647 & 4.88 / 4.87 / 0.535 \\
A3 & 1.00 & 6 / 5 / 0.438 & 1849 / 1876 / 0.780 & 0.07 / 0.28 / 0.015 \\
A3 & 0.80 & 70 / 68 / 0.425 & 139 / 96 / 0.182 & 0.22 / 0.17 / 0.843 \\
A3 & 1.50 & 1 / 1 / 1.000 & 2042 / 2048 / 0.975 & 4.89 / 4.90 / 0.427 \\
A4 & 2.00 & 1 / 0 / 0.147 & 7 / 0 / 0.087 & 0.44 / 0.78 / 0.034 \\
A4 & 0.60 & 0 / 0 / 1.000 & 0 / 0 / 1.000 & -- \\
A4 & 0.80 & 0 / 0 / 1.000 & 0 / 0 / 1.000 & -- \\
A4 & 1.00 & 0 / 0 / 1.000 & 0 / 0 / 1.000 & -- \\
A4 & 1.25 & 0 / 0 / 1.000 & 0 / 0 / 1.000 & -- \\
A4 & 1.50 & 0 / 0 / 1.000 & 0 / 0 / 1.000 & -- \\
\bottomrule
\end{tabular}

\end{table}

\end{document}
